\documentclass[12pt]{article}
\usepackage[utf8]{inputenc}
\usepackage[T1]{fontenc}
\usepackage[margin=1in]{geometry}
\usepackage{setspace}
\usepackage{xcolor}
\usepackage[colorlinks=true, urlcolor=blue!60!black, linkcolor=blue!60!black]{hyperref}
\setlength{\parindent}{2em}
\setlength{\parskip}{1.25em}
\renewcommand{\baselinestretch}{1.0}
\begin{document}
\begin{center}
{\LARGE \textit{Fulton v. City of Philadelphia}, 593 U.S. 522 (2021)\par}
\vspace{0.8em}
{\large Case Brief\par}
\end{center}
\vspace{1em}

\section{Memory Jogger}
Unanimously reversing the Third Circuit while declining to overrule \textit{Employment Division v. Smith}, the Supreme Court held that Philadelphia violated the Free Exercise Clause by refusing to contract with Catholic Social Services unless it certified same-sex foster couples, because the City's foster-care contract contained a discretionary-exemption clause that made the policy not generally applicable and unable to survive strict scrutiny (\href{https://www.courtlistener.com/opinion/4892570/fulton-v-philadelphia/}{CourtListener}; \href{https://opencasebook.org/casebooks/4501-constitutional-law/resources/12.9-fulton-v-city-of-philadelphia-2020/}{OpenCasebook}).

\section{Detailed Case Facts}
Philadelphia's foster care system relies on annual contracts between the City's Department of Human Services (DHS) and private, state-licensed foster care agencies (\href{https://opencasebook.org/casebooks/4501-constitutional-law/resources/12.9-fulton-v-city-of-philadelphia-2020/}{OpenCasebook}). Under Pennsylvania law (55 Pa. Code \S~3700.64), private agencies review and certify prospective foster families through an intensive three-to-six-month process involving background checks, medical exams, and detailed home studies assessing applicants' mental and emotional adjustment, community ties, and family relationships (\href{https://opencasebook.org/casebooks/4501-constitutional-law/resources/12.9-fulton-v-city-of-philadelphia-2020/}{OpenCasebook}). When DHS needs to place a child in its custody, it sends referrals to contracted agencies, which report available certified families (\href{https://opencasebook.org/casebooks/4501-constitutional-law/resources/12.9-fulton-v-city-of-philadelphia-2020/}{OpenCasebook}).

Catholic Social Services (CSS), a religious foster care agency affiliated with the Archdiocese of Philadelphia, had contracted with the City for over 50 years (\href{https://opencasebook.org/casebooks/4501-constitutional-law/resources/12.9-fulton-v-city-of-philadelphia-2020/}{OpenCasebook}). Guided by Catholic teaching that marriage is a sacred bond between one man and one woman, CSS viewed certifying a couple as a religious endorsement of their relationship; accordingly, it would not certify unmarried couples (regardless of sexual orientation) or married same-sex couples (\href{https://opencasebook.org/casebooks/4501-constitutional-law/resources/12.9-fulton-v-city-of-philadelphia-2020/}{OpenCasebook}). CSS did not object to certifying single gay or lesbian individuals or placing LGBTQ+ children, and no same-sex couple had ever sought certification from CSS; if one did, CSS stated it would refer the couple to one of the more than 20 other contracted agencies in Philadelphia that certified same-sex couples (\href{https://opencasebook.org/casebooks/4501-constitutional-law/resources/12.9-fulton-v-city-of-philadelphia-2020/}{OpenCasebook}).

In March 2018, after a newspaper article quoted an Archdiocesan spokesman confirming CSS's policy, the Philadelphia City Council called for an investigation into ``discrimination that occurs under the guise of religious freedom,'' and the Philadelphia Commission on Human Relations opened an inquiry (\href{https://opencasebook.org/casebooks/4501-constitutional-law/resources/12.9-fulton-v-city-of-philadelphia-2020/}{OpenCasebook}; \href{https://www.oyez.org/cases/2020/19-123}{Oyez}). During a meeting with CSS leadership, the DHS Commissioner remarked that ``things have changed since 100 years ago'' and that ``it would be great if we followed the teachings of Pope Francis, the voice of the Catholic Church'' (\href{https://opencasebook.org/casebooks/4501-constitutional-law/resources/12.9-fulton-v-city-of-philadelphia-2020/}{OpenCasebook}). Immediately afterward, DHS froze all foster child referrals to CSS and informed CSS that it would not renew its full foster care contract unless CSS agreed to certify same-sex couples (\href{https://opencasebook.org/casebooks/4501-constitutional-law/resources/12.9-fulton-v-city-of-philadelphia-2020/}{OpenCasebook}). The City asserted that CSS's refusal violated both its standard foster-care contract---specifically section 3.21 (which barred rejecting prospective foster parents based on sexual orientation ``unless an exception is granted by the Commissioner or the Commissioner's designee, in his/her sole discretion'') and section 15.1---and Philadelphia's Fair Practices Ordinance (FPO), Phila. Code \S~9-1102(1)(w), which prohibits sexual-orientation discrimination in ``public accommodations'' (\href{https://opencasebook.org/casebooks/4501-constitutional-law/resources/12.9-fulton-v-city-of-philadelphia-2020/}{OpenCasebook}; \href{https://caselaw.findlaw.com/court/us-dis-crt-ed-pen/1943877.html}{FindLaw E.D. Pa.}).

\section{Procedural History}
CSS and three affiliated foster parents---including Sharonell Fulton and Toni Lynn Simms-Busch---sued the City of Philadelphia, DHS, and the Philadelphia Commission on Human Relations under 42 U.S.C.~\S~1983 in the U.S. District Court for the Eastern District of Pennsylvania, alleging violations of the First Amendment's Free Exercise and Free Speech Clauses and seeking a temporary restraining order and preliminary injunction (\href{https://opencasebook.org/casebooks/4501-constitutional-law/resources/12.9-fulton-v-city-of-philadelphia-2020/}{OpenCasebook}; \href{https://www.supremecourt.gov/DocketPDF/19/19-123/108931/20190722174037071_Cert\%20Petition\%20FINAL.pdf}{Certiorari Petition}). Support Center for Child Advocates and Philadelphia Family Pride intervened as defendants (\href{https://www.oyez.org/cases/2020/19-123}{Oyez}).

On July 13, 2018, Judge Petrese B. Tucker denied preliminary relief (\textit{Fulton v. City of Philadelphia}, 320 F. Supp. 3d 661 (E.D. Pa. 2018)), holding that both the contractual nondiscrimination terms and the Fair Practices Ordinance applied to CSS and were neutral, generally applicable policies valid under \textit{Employment Division v. Smith}, 494 U.S. 872 (1990) (\href{https://caselaw.findlaw.com/court/us-dis-crt-ed-pen/1943877.html}{FindLaw E.D. Pa.}; \href{https://www.gwlr.org/two-surprises-in-fulton-v-city-of-philadelphia/}{GW Law Review}). On April 22, 2019, the U.S. Court of Appeals for the Third Circuit, in an opinion by Judge Thomas L. Ambro, affirmed (922 F.3d 140 (3d Cir. 2019)), concluding that Philadelphia's policies were neutral and generally applicable under \textit{Smith} and that CSS had not proven religious hostility (\href{https://www.oyez.org/cases/2020/19-123}{Oyez}; \href{https://www.gwlr.org/two-surprises-in-fulton-v-city-of-philadelphia/}{GW Law Review}).

The Supreme Court granted certiorari on February 24, 2020, including the question whether \textit{Employment Division v. Smith} should be revisited (\href{https://www.oyez.org/cases/2020/19-123}{Oyez}). Oral argument took place on November 4, 2020, and on June 17, 2021, the Supreme Court unanimously reversed and remanded (\href{https://www.oyez.org/cases/2020/19-123}{Oyez}). On remand, the District Court entered a consent permanent injunction on October 1, 2021, barring the City from enforcing the challenged condition against CSS (\href{https://clearinghouse.net/case/17588/}{Civil Rights Litigation Clearinghouse}).

\section{Judicial Votes}
\begin{itemize}
    \item \textbf{Overall Vote:} 9--0 to reverse and remand (\href{https://www.oyez.org/cases/2020/19-123}{Oyez}; \href{https://www.gwlr.org/two-surprises-in-fulton-v-city-of-philadelphia/}{GW Law Review}).
    \item \textbf{Majority Opinion (6 Justices):} Chief Justice John G. Roberts, Jr., delivered the opinion of the Court, joined by Justices Stephen G. Breyer, Sonia Sotomayor, Elena Kagan, Brett M. Kavanaugh, and Amy Coney Barrett (\href{https://www.oyez.org/cases/2020/19-123}{Oyez}).
    \item \textbf{Concurring Opinion (3 Justices):} Justice Amy Coney Barrett filed a concurring opinion joined by Justice Brett M. Kavanaugh in full and by Justice Stephen G. Breyer as to all but the first paragraph (\href{https://opencasebook.org/casebooks/4501-constitutional-law/resources/12.9-fulton-v-city-of-philadelphia-2020/}{OpenCasebook}).
    \item \textbf{Concurring in the Judgment (3 Justices):} Justice Samuel A. Alito, Jr., filed an opinion concurring in the judgment, joined by Justices Clarence Thomas and Neil M. Gorsuch (\href{https://opencasebook.org/casebooks/4501-constitutional-law/resources/12.9-fulton-v-city-of-philadelphia-2020/}{OpenCasebook}).
    \item \textbf{Concurring in the Judgment (3 Justices):} Justice Neil M. Gorsuch filed an opinion concurring in the judgment, joined by Justices Clarence Thomas and Samuel A. Alito, Jr. (\href{https://opencasebook.org/casebooks/4501-constitutional-law/resources/12.9-fulton-v-city-of-philadelphia-2020/}{OpenCasebook}).
    \item \textbf{Dissents:} None.
\end{itemize}

\section{Holding}
\begin{itemize}
    \item \textbf{Judgment:} Reversed and remanded. Philadelphia's refusal to contract with CSS for foster care services unless CSS agrees to certify same-sex couples violates the Free Exercise Clause of the First Amendment (\href{https://opencasebook.org/casebooks/4501-constitutional-law/resources/12.9-fulton-v-city-of-philadelphia-2020/}{OpenCasebook}).
    \item \textbf{Rule Announced:} Under \textit{Employment Division v. Smith}, 494 U.S. 872 (1990), a government policy is \emph{not} ``generally applicable'' if it creates a formal mechanism for individualized discretionary exemptions---regardless of whether any exemption has ever been granted (\href{https://opencasebook.org/casebooks/4501-constitutional-law/resources/12.9-fulton-v-city-of-philadelphia-2020/}{OpenCasebook}). Once a policy lacks general applicability and burdens religious exercise, it triggers strict scrutiny, requiring the government to prove a compelling interest in denying an exception to the \emph{particular} religious claimant and narrow tailoring (\href{https://opencasebook.org/casebooks/4501-constitutional-law/resources/12.9-fulton-v-city-of-philadelphia-2020/}{OpenCasebook}).
    \item \textbf{Holding vs. Dicta:} The Court held that (1) section 3.21 of Philadelphia's foster care contract is not generally applicable because it empowers the DHS Commissioner to grant exceptions in his or her ``sole discretion''; (2) Philadelphia's Fair Practices Ordinance does not apply to CSS because foster parent certification is a customized, selective assessment rather than a ``public accommodation''; and (3) the City's asserted interests fail strict scrutiny as applied to denying CSS an exception (\href{https://opencasebook.org/casebooks/4501-constitutional-law/resources/12.9-fulton-v-city-of-philadelphia-2020/}{OpenCasebook}). The Court expressly declined to decide whether \textit{Smith} should be overruled, whether Philadelphia violated the neutrality prong under \textit{Masterpiece Cakeshop}, or whether the City violated the Free Speech Clause (\href{https://opencasebook.org/casebooks/4501-constitutional-law/resources/12.9-fulton-v-city-of-philadelphia-2020/}{OpenCasebook}). All discussion of overruling \textit{Smith} and potential replacement frameworks in the Barrett, Alito, and Gorsuch concurrences is non-binding dicta.
\end{itemize}

\section{Analysis of Opinions}

\subsection{Majority Opinion (Chief Justice Roberts, joined by Justices Breyer, Sotomayor, Kagan, Kavanaugh, and Barrett)}
Chief Justice Roberts first established that Philadelphia burdened CSS's religious exercise by forcing it to choose between curtailing its historic foster-care mission and approving relationships contrary to its religious beliefs, reiterating that courts do not question the internal logic of sincere religious convictions (\href{https://opencasebook.org/casebooks/4501-constitutional-law/resources/12.9-fulton-v-city-of-philadelphia-2020/}{OpenCasebook}). Turning to \textit{Employment Division v. Smith}, he declined to reconsider \textit{Smith} because Philadelphia's policies failed \textit{Smith}'s requirement of general applicability (\href{https://opencasebook.org/casebooks/4501-constitutional-law/resources/12.9-fulton-v-city-of-philadelphia-2020/}{OpenCasebook}).

Relying on \textit{Sherbert v. Verner}, 374 U.S. 398 (1963), and \textit{Church of Lukumi Babalu Aye, Inc. v. Hialeah}, 508 U.S. 520 (1993), the Chief Justice explained that a law lacks general applicability if it invites individualized government assessment through a mechanism for discretionary exemptions (\href{https://opencasebook.org/casebooks/4501-constitutional-law/resources/12.9-fulton-v-city-of-philadelphia-2020/}{OpenCasebook}). Section 3.21 of Philadelphia's standard foster-care contract prohibited agencies from rejecting prospective foster parents based on sexual orientation ``unless an exception is granted by the Commissioner or the Commissioner's designee, in his/her sole discretion'' (\href{https://opencasebook.org/casebooks/4501-constitutional-law/resources/12.9-fulton-v-city-of-philadelphia-2020/}{OpenCasebook}). That formal exemption mechanism destroyed general applicability even though the Commissioner had never granted an exception, because its very existence invited officials to decide which reasons for noncompliance merited solicitude (\href{https://opencasebook.org/casebooks/4501-constitutional-law/resources/12.9-fulton-v-city-of-philadelphia-2020/}{OpenCasebook}). The Court also rejected the City's argument for greater leeway in managing government contractors (\href{https://opencasebook.org/casebooks/4501-constitutional-law/resources/12.9-fulton-v-city-of-philadelphia-2020/}{OpenCasebook}).

Next, the majority held that Philadelphia's Fair Practices Ordinance did not supply an independent, generally applicable bar because foster-parent certification is not a ``public accommodation'' under Phila. Code \S~9-1102(1)(w). Drawing on dictionary definitions and Pennsylvania's state antidiscrimination statute (43 Pa. Stat. Ann.~\S~954(l)), Chief Justice Roberts reasoned that public accommodations---like hotels, restaurants, and buses---are readily accessible to the general public, whereas foster certification is a selective, customized, months-long evaluation (\href{https://opencasebook.org/casebooks/4501-constitutional-law/resources/12.9-fulton-v-city-of-philadelphia-2020/}{OpenCasebook}).

Applying strict scrutiny under \textit{Gonzales v. O Centro Esp\'{i}rita Beneficente Uni\~{a}o do Vegetal}, 546 U.S. 418 (2006), the Court asked not whether Philadelphia had compelling interests in the abstract, but whether it had a compelling interest in denying an exception to CSS specifically (\href{https://opencasebook.org/casebooks/4501-constitutional-law/resources/12.9-fulton-v-city-of-philadelphia-2020/}{OpenCasebook}). Maximizing foster families would likely be \emph{helped}, not hurt, by keeping CSS in the program; liability fears were speculative because Pennsylvania delegates certification authority; and while ensuring equal treatment of gay couples is a weighty interest, the existence of a contractual exemption system fatally undermined the City's claim that its policy could brook no departures (\href{https://opencasebook.org/casebooks/4501-constitutional-law/resources/12.9-fulton-v-city-of-philadelphia-2020/}{OpenCasebook}).

\subsection{Concurring Opinion (Justice Barrett, joined by Justice Kavanaugh, and by Justice Breyer except as to Paragraph 1)}
Justice Barrett joined the Court's opinion in full but wrote separately to address \textit{Smith}. In her first paragraph (joined only by Justice Kavanaugh), she observed that while the historical record on founding-era religious exemptions is ``more silent than supportive,'' the textual and structural arguments against \textit{Smith} are compelling because it is hard to see why the Free Exercise Clause alone among First Amendment rights should offer mere protection against discrimination (\href{https://opencasebook.org/casebooks/4501-constitutional-law/resources/12.9-fulton-v-city-of-philadelphia-2020/}{OpenCasebook}).

In the remainder of her opinion (joined by both Justices Kavanaugh and Breyer), Justice Barrett expressed skepticism about simply replacing \textit{Smith}'s categorical rule with an equally categorical strict-scrutiny regime (\href{https://opencasebook.org/casebooks/4501-constitutional-law/resources/12.9-fulton-v-city-of-philadelphia-2020/}{OpenCasebook}). Overruling \textit{Smith} would raise hard doctrinal questions---such as whether religious institutions should be treated differently from individuals, whether direct and indirect burdens should receive different tiers of scrutiny, and whether pre-\textit{Smith} cases rejecting exemptions from routine laws would survive. Because strict scrutiny already governed under the contract's individualized-exemption clause and all nine Justices agreed Philadelphia failed it, she saw no reason to overrule \textit{Smith} here (\href{https://opencasebook.org/casebooks/4501-constitutional-law/resources/12.9-fulton-v-city-of-philadelphia-2020/}{OpenCasebook}).

\subsection{Concurrence in the Judgment (Justice Alito, joined by Justices Thomas and Gorsuch)}
In a sweeping 77-page concurrence, Justice Alito argued that the Court should overrule \textit{Smith} and restore \textit{Sherbert}'s strict-scrutiny standard for any law substantially burdening religious exercise (\href{https://opencasebook.org/casebooks/4501-constitutional-law/resources/12.9-fulton-v-city-of-philadelphia-2020/}{OpenCasebook}; \href{https://www.gwlr.org/two-surprises-in-fulton-v-city-of-philadelphia/}{GW Law Review}). He criticized the majority's narrow reliance on section 3.21 as a decision ``written on the dissolving paper sold in magic shops,'' warning that Philadelphia could nullify the ruling overnight simply by deleting the never-used exemption clause (\href{https://opencasebook.org/casebooks/4501-constitutional-law/resources/12.9-fulton-v-city-of-philadelphia-2020/}{OpenCasebook}).

Evaluating text, history, and \textit{stare decisis}, Justice Alito contended that the ordinary 1791 meaning of ``prohibiting the free exercise of religion'' was forbidding or hindering unrestrained religious practice, not merely banning religious discrimination (\href{https://opencasebook.org/casebooks/4501-constitutional-law/resources/12.9-fulton-v-city-of-philadelphia-2020/}{OpenCasebook}). He pointed to colonial charters and 12 of the 13 original state constitutions, which protected free exercise subject only to carveouts for ``public peace'' and ``safety''---carveouts that would be unnecessary if every neutral law automatically trumped religious practice (\href{https://opencasebook.org/casebooks/4501-constitutional-law/resources/12.9-fulton-v-city-of-philadelphia-2020/}{OpenCasebook}). He argued that \textit{Smith} was unworkable (generating confusion over ``hybrid rights,'' subjective ``targeting,'' exemptions, and COVID-19 comparators), contradicted by \textit{Hosanna-Tabor} and \textit{Masterpiece Cakeshop}, and refuted by decades of workable judicial administration under RFRA and RLUIPA (\href{https://opencasebook.org/casebooks/4501-constitutional-law/resources/12.9-fulton-v-city-of-philadelphia-2020/}{OpenCasebook}).

\subsection{Concurrence in the Judgment (Justice Gorsuch, joined by Justices Thomas and Alito)}
Justice Gorsuch sharply criticized the majority for ``dodging'' the question on which the Court granted certiorari (\href{https://opencasebook.org/casebooks/4501-constitutional-law/resources/12.9-fulton-v-city-of-philadelphia-2020/}{OpenCasebook}). He accused the majority of engaging in a ``statutory shell game'' by importing Pennsylvania's state public-accommodations law to narrow Philadelphia's broader FPO, and of executing ``a dizzying series of maneuvers'' to rewrite section 15.1 of the foster-care contract so that section 3.21's referral-stage exception swallowed the contract's blanket nondiscrimination clause (\href{https://opencasebook.org/casebooks/4501-constitutional-law/resources/12.9-fulton-v-city-of-philadelphia-2020/}{OpenCasebook}). Pointing to the ongoing litigation endured by baker Jack Phillips after \textit{Masterpiece Cakeshop}, Justice Gorsuch warned that narrow evasions force religious believers into endless litigation and argued that the Court should overrule \textit{Smith} now and work out subsidiary questions case by case (\href{https://opencasebook.org/casebooks/4501-constitutional-law/resources/12.9-fulton-v-city-of-philadelphia-2020/}{OpenCasebook}).

\section{Examples: Future Applications}
\begin{enumerate}
    \item \textbf{Same-Side (Strict Scrutiny Applies / Claimant Likely Wins): City Contracting Rule with ``Good Cause'' Waiver.} A municipal homelessness agency requires all contracted faith-based shelters to adopt a uniform staff code of conduct, ``unless the Director grants a written waiver for good cause.'' A religious shelter seeks an exemption from a provision conflicting with its religious teachings on hiring, and the Director refuses despite never having granted a waiver. Under \textit{Fulton}, the formal discretionary waiver mechanism renders the policy not generally applicable, triggering strict scrutiny even if no waiver has ever been granted (\href{https://opencasebook.org/casebooks/4501-constitutional-law/resources/12.9-fulton-v-city-of-philadelphia-2020/}{OpenCasebook}).
    \item \textbf{Same-Side (Strict Scrutiny Applies / Claimant Likely Wins): Selective Licensing Program with Secular Carveouts.} A state department of health licenses private grief-counseling nonprofits through an individualized, multi-month vetting process and prohibits licensees from declining clients on any protected basis, except when the Commissioner approves a ``programmatic specialization'' waiver. Because the selective licensing process is not a public accommodation and the waiver invites individualized discretion, refusing a religious nonprofit's accommodation triggers strict scrutiny under \textit{Fulton}.
    \item \textbf{Opposite-Side (Smith Safe Harbor Applies / Government Likely Wins): Zero-Discretion Amended Foster Contract.} Following \textit{Fulton}, a city rewrites its standard foster-care contract to eliminate section 3.21's ``sole discretion'' exception clause entirely, replacing it with an absolute, exception-free nondiscrimination rule applicable to every contracted agency and leaving no official with authority to waive it. Under \textit{Fulton}'s narrow holding---and as Justice Alito warned---eliminating the individualized-exemption mechanism restores general applicability under \textit{Smith}, insulating the contract from strict scrutiny unless the challenger proves religious hostility or secular underinclusiveness (\href{https://opencasebook.org/casebooks/4501-constitutional-law/resources/12.9-fulton-v-city-of-philadelphia-2020/}{OpenCasebook}; \href{https://www.gwlr.org/two-surprises-in-fulton-v-city-of-philadelphia/}{GW Law Review}).
    \item \textbf{Opposite-Side (Smith Safe Harbor Applies / Government Likely Wins): Exception-Free Public Accommodations Law for Commercial Venues.} A state public accommodations statute unconditionally prohibits hotels and restaurants open to the general public from refusing service based on religion or sexual orientation, with no statutory or administrative waiver mechanism. Unlike foster-parent certification, a hotel or restaurant is squarely a public accommodation, and without individualized exemptions or secular carveouts, the law remains neutral and generally applicable under \textit{Smith} (barring an independent Free Speech claim).
    \item \textbf{Fence-Sitter (Hard Case): Exception-Free Rule Coupled with Hostile Official Statements.} A city deletes all discretionary exemption clauses from its foster-care contract and amends its Fair Practices Ordinance to explicitly cover foster agencies, but during the legislative debate, several council members mock a Catholic agency's doctrine as ``bigotry hiding behind outdated dogma.'' Even though the revised policy satisfies \textit{Fulton}'s general-applicability test, it falls into the unresolved boundary of \textit{Masterpiece Cakeshop} and \textit{Lukumi}: courts must decide whether legislators' hostile remarks (as opposed to adjudicators' remarks) strip a facially neutral, exception-free law of neutrality under \textit{Smith} (\href{https://opencasebook.org/casebooks/4501-constitutional-law/resources/12.9-fulton-v-city-of-philadelphia-2020/}{OpenCasebook}).
\end{enumerate}

\section{Critique}
\begin{itemize}
    \item \textbf{Originalist and Textualist Critiques:} Conservative originalists and textualists, echoing Justices Alito and Gorsuch, criticize \textit{Fulton} as a missed opportunity that ``kicked the can down the road'' (\href{https://red.library.usd.edu/cgi/viewcontent.cgi?article=1102&context=sdlrev}{South Dakota Law Review}). They argue that \textit{Smith} cannot be reconciled with the original public meaning of ``prohibiting the free exercise of religion'' and that Chief Justice Roberts's contorted reading of Philadelphia's contract and Fair Practices Ordinance compromised textual fidelity simply to manufacture unanimity and avoid overruling \textit{Smith} (\href{https://opencasebook.org/casebooks/4501-constitutional-law/resources/12.9-fulton-v-city-of-philadelphia-2020/}{OpenCasebook}; \href{https://con.law/print-articles/dwyer-smiths-last-stand-foster-care}{ConLaw}). At the same time, some originalist scholars highlight Justice Barrett's concurrence as a disciplined methodological check: because founding-era history is genuinely contested and silent on judicial mandates for religious exemptions, rushing to replace \textit{Smith} with universal strict scrutiny risks unprincipled judicial balancing (\href{https://www.gwlr.org/two-surprises-in-fulton-v-city-of-philadelphia/}{GW Law Review}; \href{https://www.scotusblog.com/2021/06/court-requires-religious-exemption-but-leaves-many-questions-unanswered/}{SCOTUSblog}).
    \item \textbf{Progressive and Doctrinal Critiques:} Progressive and church-state scholars such as Ira C. Lupu, Robert W. Tuttle, and Nelson Tebbe criticize the majority for stretching the obscure ``individualized exemptions'' doctrine beyond recognition (\href{https://www.gwlr.org/two-surprises-in-fulton-v-city-of-philadelphia/}{GW Law Review}; \href{https://michiganlawreview.org/journal/the-politics-of-proportionality/}{Michigan Law Review}). In their view, treating a \emph{never-exercised} boilerplate discretion clause as fatal to general applicability---and then using that unused clause to declare that the City's antidiscrimination interest cannot survive strict scrutiny---hollows out \textit{Smith} without formally overruling it (\href{https://www.gwlr.org/two-surprises-in-fulton-v-city-of-philadelphia/}{GW Law Review}; \href{https://www.acslaw.org/analysis/acs-journal/2020-2021-acs-supreme-court-review/the-radical-uncertainty-of-free-exercise-principles-a-comment-on-fulton-v-city-of-philadelphia/}{ACS Supreme Court Review}). Critics also argue that the Court ignored the dignitary and practical harms to LGBTQ+ prospective foster parents and foster youth when a taxpayer-funded, government-delegated public service turns away qualified couples, and faulted the Court's minimalism for evading an honest proportionality inquiry between competing constitutional commitments (\href{https://michiganlawreview.org/journal/the-politics-of-proportionality/}{Michigan Law Review}; \href{https://www.cardozolawreview.com/remembering-who-foster-care-is-for-public-accommodation-and-other-misconceptions-and-missed-opportunities-in-fulton-v-city-of-philadelphia/}{Cardozo Law Review}).
    \item \textbf{Underlying Competing Values:} \textit{Fulton} sits at the sharp intersection of \emph{religious liberty and pluralism} (enabling faith-based charities to participate in public life without violating core theological tenets) and \emph{civic equality and non-discrimination} (ensuring that publicly funded government services treat LGBTQ+ citizens with equal dignity and access, without forcing the state to subsidize exclusion).
\end{itemize}

\section{Key Quotations}

\begin{quote}
``A law is not generally applicable if it `invite[s]' the government to consider the particular reasons for a person's conduct by providing `\,`a mechanism for individualized exemptions.'\,'''
\end{quote}
\noindent--- Chief Justice John G. Roberts, Jr., Opinion of the Court, \textit{Fulton v. City of Philadelphia}, 593 U.S. 522 (2021) (\href{https://opencasebook.org/casebooks/4501-constitutional-law/resources/12.9-fulton-v-city-of-philadelphia-2020/}{OpenCasebook}; \href{https://www.courtlistener.com/opinion/4892570/fulton-v-philadelphia/}{CourtListener}).

\begin{quote}
``The creation of a formal mechanism for granting exceptions renders a policy not generally applicable, regardless whether any exceptions have been given, because it `invite[s]' the government to decide which reasons for not complying with the policy are worthy of solicitude, \textit{Smith}, 494 U.S. at 884---here, at the Commissioner's `sole discretion.' ''
\end{quote}
\noindent--- Chief Justice John G. Roberts, Jr., Opinion of the Court, \textit{Fulton v. City of Philadelphia}, 593 U.S. 522 (2021) (\href{https://opencasebook.org/casebooks/4501-constitutional-law/resources/12.9-fulton-v-city-of-philadelphia-2020/}{OpenCasebook}).

\begin{quote}
``Rather than rely on `broadly formulated interests,' courts must `scrutinize[] the asserted harm of granting specific exemptions to particular religious claimants.' ''
\end{quote}
\noindent--- Chief Justice John G. Roberts, Jr., Opinion of the Court, \textit{Fulton v. City of Philadelphia}, 593 U.S. 522 (2021) (\href{https://opencasebook.org/casebooks/4501-constitutional-law/resources/12.9-fulton-v-city-of-philadelphia-2020/}{OpenCasebook}).

\begin{quote}
``As a matter of text and structure, it is difficult to see why the Free Exercise Clause---lone among the First Amendment freedoms---offers nothing more than protection from discrimination.''
\end{quote}
\noindent--- Justice Amy Coney Barrett, Concurring Opinion, \textit{Fulton v. City of Philadelphia}, 593 U.S. 522 (2021) (\href{https://opencasebook.org/casebooks/4501-constitutional-law/resources/12.9-fulton-v-city-of-philadelphia-2020/}{OpenCasebook}).

\begin{quote}
``This decision might as well be written on the dissolving paper sold in magic shops. The City has been adamant about pressuring CSS to give in, and if the City wants to get around today's decision, it can simply eliminate the never-used exemption power.''
\end{quote}
\noindent--- Justice Samuel A. Alito, Jr., Concurring in the Judgment, \textit{Fulton v. City of Philadelphia}, 593 U.S. 522 (2021) (\href{https://opencasebook.org/casebooks/4501-constitutional-law/resources/12.9-fulton-v-city-of-philadelphia-2020/}{OpenCasebook}).

\end{document}
